\BourbakiExercises{Exercises}

\begin{enumerate}
\def\labelenumi{\arabic{enumi})}
\tightlist
\item
  \begin{enumerate}
  \def\labelenumii{\alph{enumii})}
  \tightlist
  \item
    Give an example of a Noetherian module M such that every nonzero endomorphism of M is injective and that there exist nonzero nonbijective endomorphisms of M.
  \end{enumerate}
\end{enumerate}

\begin{enumerate}
\def\labelenumi{\alph{enumi})}
\setcounter{enumi}{1}
\tightlist
\item
  Give an example of an Artinian module M such that every nonzero endomorphism of M is surjective and that there exist nonzero nonbijective endomorphisms of M.
\end{enumerate}

\begin{enumerate}
\def\labelenumi{\arabic{enumi})}
\setcounter{enumi}{1}
\item
  Let \(\mathrm{A}\) be a ring in which every increasing sequence \((\mathfrak{a}_n)_{n\in \mathbf{N}}\) of two-sided ideals is stationary (for example, a left or right Noetherian ring). Prove that every surjective endomorphism of the ring \(\mathrm{A}\) is bijective.
\item
  Give an example of a ring that has two elements u, v such that uv = 1 and \(vu \neq 1\) . Prove that such a ring is not the union of a directed family of Noetherian subrings.
\item
  Let M be an A-module and p and q two projectors of M.
\end{enumerate}

\begin{enumerate}
\def\labelenumi{\alph{enumi})}
\item
  We have \(p(\mathrm{M}) \subset q(\mathrm{M})\) if and only if we have qp = p; the kernel \(\operatorname{Ker}(p)\) is contained in \(\operatorname{Ker}(q)\) if and only if we have qp = q.
\item
  The sum \(p + q\) is a projector if and only if we have pq = -qp. If \(p + q\) is a projector and if in M, the relation 2x = 0 implies x = 0, then we have pq = 0 and qp = 0.
\item
  Prove that the following properties are equivalent:
\end{enumerate}

\begin{enumerate}
\def\labelenumi{(\roman{enumi})}
\item
  We have \(pq = qp\) .
\item
  We have \(q(p(\mathbb{M})) \subset p(\mathbb{M})\) and \(q(\operatorname{Ker}(p)) \subset \operatorname{Ker}(p)\) .
\item
  There exists a sequence \((p_1, p_2, p_3, p_4)\) of pairwise orthogonal projectors of M such that \(p = p_1 + p_2\) , \(q = p_1 + p_3\) , and \(1_{\mathrm{M}} = p_1 + p_2 + p_3 + p_4\) .
\end{enumerate}

Such a sequence, when it exists, is unique: we have \(p_{1} = pq\) .

\begin{enumerate}
\def\labelenumi{\arabic{enumi})}
\setcounter{enumi}{4}
\tightlist
\item
  Let \(\mathrm{A}\) be a ring.
\end{enumerate}

\begin{enumerate}
\def\labelenumi{\alph{enumi})}
\item
  An element e of A is idempotent if and only if the mapping \(x \mapsto xe\) is a projector of the left A-module \(A_{s}\) (resp. the mapping \(x \mapsto ex\) is a projector of the right A-module \(A_{d}\) ).
\item
  Let \(e_{1}, e_{2}\) be idempotents in the ring A. Prove that the group \(\mathrm{Hom}_{\mathrm{A}}(\mathrm{Ae}_{1}, \mathrm{Ae}_{2})\) is isomorphic to \(e_{1}A \cap Ae_{2} = e_{1}Ae_{2}\) and that the ring \(\mathrm{End}_{\mathrm{A}}(\mathrm{Ae}_{1})\) is isomorphic to the ring \(e_{1}Ae_{1}\) .
\item
  Let \(e_1, e_2\) be idempotents in A. Prove that the following properties are equivalent:
\end{enumerate}

\begin{enumerate}
\def\labelenumi{(\roman{enumi})}
\item
  \(Ae_{1} = Ae_{2}.\)
\item
  \(e_{1}e_{2}=e_{1}\) and \(e_{2}e_{1}=e_{2}\) .
\item
  \((1 - e_{1})\mathrm{A} = (1 - e_{2})\mathrm{A}.\)
\end{enumerate}

When they are satisfied, there exists an invertible element \(a \in \mathrm{A}\) such that \(e_2 = ae_1a^{-1}\) .

\begin{enumerate}
\def\labelenumi{\alph{enumi})}
\setcounter{enumi}{3}
\tightlist
\item
  Let \(e_{1}, e_{2}\) be two idempotents in A. Prove that the following properties are equivalent:
\end{enumerate}

\begin{enumerate}
\def\labelenumi{(\roman{enumi})}
\item
  The left A-modules \(Ae_{1}\) and \(Ae_{2}\) are isomorphic.
\item
  The right A-modules \(e_1\mathrm{A}\) and \(e_2\mathrm{A}\) are isomorphic.
\item
  There exist elements \(x, y\) of \(A\) such that \(xy = e_1\) and \(yx = e_2\) .
\end{enumerate}

When they are satisfied, we say that \(e_{1}\) and \(e_{2}\) are equivalent idempotents.

\begin{enumerate}
\def\labelenumi{\alph{enumi})}
\setcounter{enumi}{4}
\tightlist
\item
  If idempotents in A are equivalent and belong to the center of A, then they are equal.
\end{enumerate}

\(f\) ) Let \((e_i)_{i\in \mathrm{I}}\) and \((f_{i})_{i\in \mathrm{I}}\) be finite families of idempotents in A such that \(e_i e_j = 0\) and \(f_{i} f_{j} = 0\) for \(i\in \mathrm{I}, j\in \mathrm{I}, i\neq j\) and that \(e_i\) is equivalent to \(f_{i}\) for every \(i\in \mathrm{I}\) . Then \(e = \sum_{i\in \mathrm{I}}e_i\) and \(f = \sum_{i\in \mathrm{I}}f_{i}\) are equivalent idempotents in A.

¶ 6) Let A be a ring and u and v two elements of A such that uv = 1 and vu ≠ 1. For i ≥ 1 and j ≥ 1, we set \(e_{ij} = v^{i-1}u^{j-1} - v^i u^j\) .

\begin{enumerate}
\def\labelenumi{\alph{enumi})}
\item
  Prove that we have \(e_{ij} \neq 0\) and \(e_{ij}e_{hk} = \delta_{jh}e_{ik}\) for \(i,j,k,h\) integers \(\geqslant 1\) .
\item
  Prove that the left ideals \(Ae_{ii}\) of A (for \(i \geqslant 1\) ) are nonzero, pairwise isomorphic as A-modules (use Exercise 5 d)), and that their sum is direct.
\item
  Prove that we have \(u(v + e_{1i}) = 1\) for every \(i \geqslant 1\) and \(e_{1i} \neq e_{1j}\) for \(i \geqslant 1, j \geqslant 1, i \neq j\) .
\item
  Let B be a ring that has the equivalent properties of Exercise 7 of VIII, p.~15, a). Deduce from the above that every left or right invertible element of B is invertible.
\end{enumerate}

\begin{enumerate}
\def\labelenumi{\arabic{enumi})}
\setcounter{enumi}{6}
\tightlist
\item
  Denote by A (resp. a) the set of rational fractions of the form \(\mathrm{P}/(\mathrm{X}^{2}+1)^{n}\) , where n is an integer \(\geqslant0\) and \(P\in \mathbf{Q}[X]\) is a polynomial of degree \(\leqslant2n\) (resp. \(\leqslant2n-1\) ). a) Prove that A is a subring of the field \(\mathbf{Q}(X)\) and that a is an ideal of the ring A. b) The A-modules \(a\oplus a\) and \(A\oplus A\) are isomorphic. The A-module a is projective and finitely generated but is not isomorphic to A.
\end{enumerate}

\begin{enumerate}
\def\labelenumi{\alph{enumi})}
\setcounter{enumi}{2}
\tightlist
\item
  There exists a quadratic extension \(\mathbf{K}\) of the field \(\mathbf{Q}\) such that the \(\mathrm{A}_{(\mathrm{K})}\) -modules \(\mathfrak{a}_{(\mathrm{K})}\) and \(\mathrm{A}_{(\mathrm{K})}\) are isomorphic.
\end{enumerate}

\begin{enumerate}
\def\labelenumi{\arabic{enumi})}
\setcounter{enumi}{7}
\item
  Let A be a left Artinian ring, and let M and N be two supplementary submodules in the A-module \(A_{s}^{n}\) . If M is a free A-module, then N is a free A-module.
\item
  Let A be a complete topological ring and m a two-sided ideal of A. Suppose that the ring A/m is a field, that every element of m is topologically nilpotent (Gen.~Top., III, §6, p.~317, Exercise 9), and that there exists a fundamental system of neighborhoods of 0 in A consisting of subgroups of the additive group of A. Then the ring A is local and m is its maximal ideal. In particular, for every prime p, the ring \(\mathbf{Z}_{p}\) of p-adic integers (V, §12, No.~3, p.~96) is local, and its maximal ideal is \(p\mathbf{Z}_{p}\) .
\end{enumerate}

¶ 10) Let p be a prime number, q a power of p, G a finite group of order q, and A a commutative ring. Let \((e_{g})\) be the canonical basis of the algebra A{[}G{]} of the group G over A.

\begin{enumerate}
\def\labelenumi{\alph{enumi})}
\item
  The set \(\mathfrak{a}\) of elements \(\sum a_{g}e_{g}\) of A{[}G{]} such that \(\sum a_{g} = 0\) is a two-sided ideal of A{[}G{]}. If the ring A is of characteristic \(p\) (V, §1, No.~2, p.~2), then we have \(x^{q} = 0\) for every \(x\in \mathfrak{a}\) (use induction on \(q\) ; if \(q\neq 1\) , choose an element \(h\) of order \(p\) of the center of G (I, §6, No.~5, p.~77, Corollary of Proposition 11) and deduce from the induction hypothesis that there exists a \(y\in \mathrm{A}[\mathrm{G}]\) such that \(x^{q / p} = (1 - e_h)y\) ).
\item
  Suppose that A is a local ring, with maximal ideal m, and that the field A/m is of characteristic p.~Prove that the ring A{[}G{]} is local and that its maximal ideal is the set of elements \(\sum a_{g}e_{g}\in A[G]\) such that \(\sum a_{g}\in m\) .
\end{enumerate}

(Let S be a simple A{[}G{]}-module; prove that S is annihilated by m and then, by applying Proposition 11 of I, §6, No.~5, p.~76, to the Z{[}G{]}-submodule generated by a nonzero element of S, that S contains an element invariant under G.)

\begin{enumerate}
\def\labelenumi{\arabic{enumi})}
\setcounter{enumi}{10}
\tightlist
\item
  Let p be a prime number, G a group that does not contain any elements of order p, A a commutative ring of characteristic p that does not contain any nonzero nilpotent elements, A{[}G{]} the algebra of G over A, and \((e_{g})\) the canonical basis of A{[}G{]} over A. Prove that the ring A{[}G{]} has no nonzero nil ideal (observe that if an element \(\sum a_{g}e_{g} \in A[G]\) is nilpotent, then we have \(a_{1} = 0\) ).
\end{enumerate}

¶ 12) Let K be a commutative field, E and F extensions of K, and A the ring \(E \otimes_{K} F\) .

\begin{enumerate}
\def\labelenumi{\alph{enumi})}
\tightlist
\item
  If E and F are transcendental extensions of K, then the ring A is not local. b) Suppose that E is an algebraic extension of K, and denote the separable closures of K in E and F, respectively, by \(E_{s}\) and \(F_{s}\) . The following properties are equivalent:
\end{enumerate}

\begin{enumerate}
\def\labelenumi{(\roman{enumi})}
\item
  The ring A is local.
\item
  The ring \(\mathrm{E}\otimes_{\mathrm{K}}\mathrm{F}_s\) is an integral domain.
\item
  The ring \(E \otimes_{K} F_{s}\) is a field.
\item
  The ring \(E_{s} \otimes_{K} F_{s}\) is a field.
\end{enumerate}

When they hold, the maximal ideal \(\mathfrak{m}\) of A is the set of nilpotent elements of A, and every composite extension of E and F is isomorphic to \(A / \mathfrak{m}\) .

\begin{enumerate}
\def\labelenumi{\arabic{enumi})}
\setcounter{enumi}{12}
\item
  Let A be a local integral domain that is not a field (for example \(\mathbf{Z}_{(p)}\) , cf.~VIII, p.~26, Example 5), and let a be a nonzero noninvertible element of A. Denote the A-module \(\mathrm{A}^{(\mathbf{N})}\) by M and its canonical basis by \((e_{n})_{n\in\mathbf{N}}\) . Prove that the sequence \((f_{n})_{n\in\mathbf{N}}\) , where \(f_{2m}=e_{2m}\) and \(f_{2m+1}=e_{2m+1}+ae_{2m+2}\) for every \(m\geqslant0\) , is a basis of M over A. Prove that the submodule \(N=\sum_{m\geqslant0}\mathrm{A}(e_{2m}+ae_{2m+1})\) of M is a direct factor but that there is no subset J of N such that \(\sum_{n\in J}Af_{j}\) is supplementary to N. Deduce that in Corollary 5 of VIII, p.~35, we cannot leave out the assumption that I is finite.
\item
  Let A be a ring, M an A-module, and N a primordial direct factor of M. Let \(M = \bigoplus_{i \in I} M_i\) be a decomposition of M into a direct sum, and let \((p_i)_{i \in I}\) be the associated family of projectors. Prove that there exists an \(i \in I\) such that \(p_i\) induces an isomorphism from N to a direct factor of \(M_i\) .
\item
  Let A be a ring, and let M, N, P be A-modules. If the modules \(M \oplus P\) and \(N \oplus P\) are isomorphic and the module P is primordial, then the modules M and N are isomorphic.
\item
  Let A be a ring, M an A-module of finite length, and \((\mathrm{M}_{i})_{i\in\mathrm{I}}\) a family of indecomposable submodules of M with direct sum M. Let L be an indecomposable A-module of finite length and \(I_{L}\) the set of indices \(i\in I\) such that \(M_{i}\) is isomorphic to L. Use an example to show that the submodule \(\sum_{i\in I_{L}}M_{i}\) of M is not necessarily stable under the automorphisms of M.
\item
  Let \(\mathrm{A}\) be a ring, \(n\) an integer, and \(X \in \mathbf{GL}_n(\mathrm{A})\) .
\end{enumerate}

\begin{enumerate}
\def\labelenumi{\alph{enumi})}
\item
  Suppose that the ring A is local. Prove that there exist matrices P, Q in \(\mathbf{GL}_{n}(\mathrm{A})\) that are products of matrices of the form \(B_{ij}(\lambda)\) (II, §10, No.~13, p.~361) and an invertible element \(\delta\) of A, such that we have \(PXQ = \text{diag}(1, \ldots, 1, \delta)\) (follow the proof of Proposition 14 of loc. cit.).
\item
  Suppose that the ring A is Euclidean (VII, §1, p.~49, Exercise 7). Prove the same result.
\item
  If \(\mathrm{A}\) is a local ring or a Euclidean commutative ring, then the group \(\mathbf{S}\mathbf{L}_n(\mathbf{A})\) is generated by the matrices \(B_{ij}(\lambda)\) .
\item
  Suppose that A is local, and denote its maximal ideal by m. Prove that every matrix in \(\mathbf{M}_{n}(\mathrm{A})\) with invertible image in \(\mathbf{M}_{n}(\mathrm{A}/\mathfrak{m})\) is itself invertible (use Lemma 1 of II, §10, No.~13, p.~361).
\end{enumerate}

\begin{enumerate}
\def\labelenumi{\arabic{enumi})}
\setcounter{enumi}{17}
\tightlist
\item
  Let A be a local ring, E a free A-module, and x a nonzero element of E. Let \((e_{i})_{i\in I}\) be a basis of E in which the number of nonzero coordinates of x is the least possible. Let \(x=\sum a_{i}e_{i}\) be the representation of x in this basis; denote the set of indices \(i\in I\) such that \(a_{i}\neq0\) by J.
\end{enumerate}

\begin{enumerate}
\def\labelenumi{\alph{enumi})}
\tightlist
\item
  Prove that none of the \(a_{i}\) belong to the right ideal of A generated by the others.
\item
  Let F and G be supplementary submodules of E such that \(x \in F\) ; for \(i \in I\) , write \(e_{i} = f_{i} + g_{i}\) with \(f_{i} \in F\) and \(g_{i} \in G\) . Prove that the family consisting of the \(f_{j}\) for \(j \in J\) and the \(e_{i}\) for \(i \in I - J\) is a basis of E (if \(g_{j} = \sum_{k \in I} b_{jk} e_{k}\) , observe that the coordinates \(b_{jk}\) for j, k in J belong to the maximal ideal of A, and apply Exercise 17 d)). Deduce that there exists a free submodule of F that contains x and is a direct factor of F.
\end{enumerate}

\(c\) ) Prove that every projective A-module is free (use Kaplansky's theorem (II, §2, p.~386, Exercise 2) to reduce to the case of a module with a countable family of generators, and apply \(b\) )).

\begin{enumerate}
\def\labelenumi{\arabic{enumi})}
\setcounter{enumi}{18}
\item
  Let A be a commutative ring, M a finitely generated A-module, u an automorphism of M, and N a submodule of M that is stable under u. Prove that we have \(u(N) = N\) (use Corollary 1, d) of VIII, p.~28).
\item
  Let A be ring, let M be an A-module of finite length, and let N be an A-module. Let m and n be strictly positive integers such that \(M^{m}\) and \(N^{n}\) are isomorphic. Prove that there exist an A-module P and integers p, q such that M is isomorphic to \(P^{p}\) , N is isomorphic to \(P^{q}\) , and mp = nq.
\end{enumerate}

